\documentclass[12pt]{article}

\usepackage{geometry}
\usepackage{longtable}
\usepackage{hyperref}
\usepackage{array}

\geometry{margin=1in}

\title{
    \textbf{Snapshot 2 Objective: Checkpoint 1} \\
    JPL Lunar Detection Pipeline
}
\author{
    Software Engineering Tools Final Project \\
    Course: CS3338 \\
    Team Members: [Name 1], [Name 2], [Name 3], [Name 4]
}
\date{[Insert Date]}

\begin{document}

\maketitle
\newpage

\section{Objective}

The objective of Snapshot 2 is to add the first major feature to the JPL Lunar Detection Pipeline. For this checkpoint, the main new feature will be simulated boulder detection.

This snapshot also introduces TestRail documentation and updates the project documents based on the new feature.

\section{Reflection on Snapshot 1}

During Snapshot 1, the team created the foundation of the project. The main completed items were:

\begin{itemize}
    \item GitHub repository structure.
    \item Initial README file.
    \item Software Design Document.
    \item Software Requirements Specification.
    \item Design Specification.
    \item User Manual.
    \item Snapshot Objective documents.
    \item Initial Jira task planning.
    \item Initial mock data planning.
\end{itemize}

The project is now ready to add a new simulated feature and begin testing documentation.

\section{New Feature: Boulder Detection}

The new feature for Snapshot 2 is simulated boulder detection. The system will use mock data to represent boulders found in lunar surface images.

\subsection{Feature Description}

Boulder detection will identify rock-like objects on the lunar surface. In a real mission, boulders may be important because they can create navigation hazards for rovers or landing systems.

For this project, boulder detection will be simulated using JSON mock data.

\subsection{Example Boulder Detection Data}

\begin{verbatim}
{
  "detection_id": "B-001",
  "image_id": "IMG-001",
  "feature_type": "Boulder",
  "confidence": 0.87,
  "estimated_size_meters": 3.2,
  "bounding_box": {
    "x": 320,
    "y": 180,
    "width": 25,
    "height": 22
  },
  "risk_level": "Medium"
}
\end{verbatim}

\section{Goals for Snapshot 2}

The goals for Snapshot 2 are:

\begin{itemize}
    \item Add boulder detection mock data.
    \item Update the SDD with the boulder detection feature.
    \item Update the SRS with new boulder detection requirements.
    \item Update the Design Specification.
    \item Update the User Manual.
    \item Add Jira tasks for boulder detection.
    \item Add Jira bug examples related to boulder detection.
    \item Create the first TestRail test cases.
    \item Export or save the Snapshot 2 TestRail summary report.
    \item Update the workflow diagram.
\end{itemize}

\section{Planned Jira Tasks}

Example Jira tasks for Snapshot 2 include:

\begin{longtable}{|p{1.5in}|p{4.5in}|}
\hline
\textbf{Issue Type} & \textbf{Task} \\
\hline
Task & Add boulder detection mock data file. \\
\hline
Task & Update SDD with boulder detection component. \\
\hline
Task & Update SRS with boulder detection requirements. \\
\hline
Task & Update User Manual with boulder detection instructions. \\
\hline
Task & Create TestRail test cases for boulder detection. \\
\hline
Bug & Boulder confidence score displays as decimal instead of percentage. \\
\hline
Bug & Boulder risk level does not appear on the Results Page. \\
\hline
\end{longtable}

\section{TestRail Plan}

Snapshot 2 will include the first TestRail report.

Example test cases include:

\begin{longtable}{|p{1.2in}|p{2.4in}|p{2.4in}|}
\hline
\textbf{Test ID} & \textbf{Test Case} & \textbf{Expected Result} \\
\hline
TC-001 & Upload valid lunar image & Image is accepted and processing starts. \\
\hline
TC-002 & Upload invalid file type & System rejects file and displays error. \\
\hline
TC-003 & View boulder detections & Boulder results appear on Results Page. \\
\hline
TC-004 & Check boulder confidence score & Confidence score displays correctly as a percentage. \\
\hline
TC-005 & Check boulder risk level & Risk level appears for each boulder detection. \\
\hline
\end{longtable}

\section{Documents to Update}

The following documents should be updated during Snapshot 2:

\begin{itemize}
    \item SDD.tex
    \item SRS.tex
    \item DesignSpec.tex
    \item UserManual.tex
    \item README.md
    \item Snapshot2\_Checkpoint1.tex
\end{itemize}

\section{Expected Deliverables}

By the end of Snapshot 2, the project should include:

\begin{itemize}
    \item Boulder detection mock data.
    \item Updated SDD.
    \item Updated SRS.
    \item Updated Design Specification.
    \item Updated User Manual.
    \item Updated README.
    \item Jira tasks and bugs for Snapshot 2.
    \item Snapshot 2 TestRail report.
    \item Updated workflow diagram.
\end{itemize}

\section{Reflection Placeholder}

This section will be completed after Snapshot 2 is finished. The team will describe what was completed, what issues came up, and what should be improved for Snapshot 3.

\end{document}